\section{Summary and Research Gaps} \label{sec:28-background-summary}

\begin{foldable}
  This chapter has built the technical foundation for the thesis along two parallel tracks.
\end{foldable}

\begin{foldable}
  Track~A (Sections~\ref{sec:21-background-deeplearning}--\ref{sec:25-background-synthetic}) established the image knowledge distillation pipeline: the deep learning architectures whose deployment cost motivates compression;
  the compression landscape that positions KD as the preferred method; the standard KD framework and the rich soft-target signal it transfers; the increasingly restricted settings --- few-shot, data-free white-box, data-free black-box --- that arise in practice;
  and the generative and contrastive tools that must be enlisted when the standard assumptions fail.
\end{foldable}

\begin{foldable}
  Track~B (Sections~\ref{sec:26-background-datasetdistillation}--\ref{sec:27-background-importanceselection}) established the text dataset distillation pipeline: the language modelling and \ac{NLP} task foundations;
  the dataset distillation objective and why image methods do not transfer to text;
  the prior text DD methods and their shared limitations;
  and the importance-scoring and optimal-transport machinery that provides principled, coverage-aware selection.
\end{foldable}

\begin{foldable}
  The unifying thread is distributional fidelity.
  In both tracks, the compact artefact --- a student model or a surrogate corpus --- must cover the distribution of the teacher or full dataset.
  Coverage failures manifest differently in the two tracks (mode collapse in GAN-based synthesis;
  coverage gaps in clustering-based corpus selection) but stem from the same root cause: the optimization objective used to produce the artefact does not directly penalise distributional divergence from the target.
\end{foldable}

\begin{foldable}
  It is worth being precise about why this indirect penalization matters.
  Existing methods do not ignore distributional fidelity; they pursue it through proxy objectives that correlate with distributional coverage without measuring it directly.
  \ac{GAN} training minimises a divergence between the generated and real distributions, but the practical objective --- adversarial cross-entropy or Wasserstein distance estimated on a finite batch --- is a tractable surrogate, not a direct measure of how well the generated distribution covers the target.
  Gradient-matching objectives in dataset distillation enforce that a surrogate corpus induces similar training dynamics to the full dataset, which correlates with coverage but does not bound it.
  Clustering-based selection methods use intra-cluster distance as a proxy for representativeness, which can be dominated by high-density regions and systematically under-select low-frequency but semantically distinct examples.
  This indirection is not a failure of prior design choices --- direct minimization of distributional divergence between a generated or selected artefact and the target distribution is computationally expensive, typically requiring approximations whose cost scales poorly with dimension and sample size.
  The tractable proxies are therefore principled engineering decisions, not oversights.
  What the thesis argues is that advances in contrastive representation learning, importance weighting, and discrete optimal transport now make direct penalization tractable in the three specific settings under study.
  Each of the three contributions makes this direct minimization feasible in its respective context; without this framing, the three methods appear as independent engineering improvements rather than as instances of a single principle applied in different regimes.
\end{foldable}

\begin{foldable}
  Table~\ref{tab:research-gaps} summarises the three settings this thesis addresses, the representative prior methods, their core limitations, and the corresponding contributions.
\end{foldable}

\begin{table}[t]
  \centering
  \caption{Research gaps addressed by the three thesis contributions.}
  \label{tab:research-gaps}
  \begin{tabular}{p{3.2cm}p{3.2cm}p{3.8cm}p{2.8cm}}
    \hline
    \textbf{Setting} & \textbf{Representative prior work} & \textbf{Core limitation} & \textbf{Thesis contribution} \\
    \hline
    Black-box few-shot KD &
    BBKD, FS-BBT &
    GAN mode collapse; no explicit diversity mechanism &
    DivBFKD (Ch.~\ref{ch:30-research01}) \\[4pt]
    Black-box data-free KD &
    ZSDB3~\cite{wang2021zero}, IDEAL~\cite{zhang2022ideal}, DFHL-RS~\cite{yuan2024data} &
    Poor domain coverage; hard-label information bottleneck &
    DIPKD (Ch.~\ref{ch:40-research02}) \\[4pt]
    Text dataset distillation &
    DiLM~\cite{maekawa2024dilm}, DaLLME~\cite{tao2024textual} &
    Hard-sample bias; no distributional coverage guarantee &
    TAKE (Ch.~\ref{ch:50-research03}) \\
    \hline
  \end{tabular}
\end{table}

\begin{foldable}
  Three open problems follow directly from the analysis.
\end{foldable}

\begin{foldable}
  \textbf{Gap~1 $\rightarrow$ Chapter~\ref{ch:30-research01}.}
  In black-box few-shot knowledge distillation, how can GAN-based synthesis be directed to produce data that is both high-fidelity and semantically diverse, without access to the teacher's internal representations, when the generator is itself trained on a few-shot dataset that already lacks diversity?
\end{foldable}

\begin{foldable}
  \textbf{Gap~2 $\rightarrow$ Chapter~\ref{ch:40-research02}.}
  In black-box data-free knowledge distillation, how can a student simultaneously generate synthetic data that covers the teacher's full input distribution and recover a soft distillation signal approximating the teacher's class probabilities --- both without any real training data and with access to only hard-label teacher queries?
\end{foldable}

\begin{foldable}
  \textbf{Gap~3 $\rightarrow$ Chapter~\ref{ch:50-research03}.}
  In text dataset distillation, how can per-sample importance scores be corrected for hard-sample bias, and how can the selected corpus be formally guaranteed to cover the importance-reweighted source distribution?
\end{foldable}

\begin{foldable}
  Chapters~\ref{ch:30-research01}, \ref{ch:40-research02}, and~\ref{ch:50-research03} each address one of these gaps.
  Chapter~\ref{ch:70-conclusion} revisits them collectively and identifies the unifying principle that connects all three solutions: that explicitly minimising distributional divergence --- whether through contrastive diversity enforcement, importance-aware generation, or optimal transport selection --- is the decisive factor separating the proposed methods from their predecessors.
\end{foldable}
